\section{The EviMatch System}\label{sec:method}
\begin{figure*}[t]
\centering
% System architecture (TikZ). Included from main.tex inside figure*.
\begin{tikzpicture}[
  font=\footnotesize,
  box/.style={draw, rounded corners=2pt, minimum height=8mm, align=center, inner sep=3pt, fill=#1},
  box/.default=white,
  arr/.style={-{Latex[length=2mm]}, thick},
  lab/.style={font=\scriptsize\itshape, text=black!70}
]
\node[box=black!6] (res) {Resume\\(PDF/DOCX/TXT)};
\node[box=black!6, below=5mm of res] (jd) {Job\\description};
\node[box=blue!10, right=9mm of res] (red) {Blind\\redaction};
\node[box=blue!10, right=9mm of jd] (req) {Requirement\\parser};
\node[box=blue!10, right=9mm of red] (sent) {Sentence\\units};
\node[box=blue!10, right=9mm of req] (sk) {Skill\\lexicon};
\node[box=orange!15, right=9mm of sent, yshift=-6mm] (feat) {Pair features\\lexical $\cdot$ semantic\\skills $\cdot$ experience};
\node[box=orange!15, right=9mm of feat] (head) {Fit model\\(prior-corrected)};
\node[box=green!12, right=9mm of head, yshift=7mm] (out) {Score\\+ band};
\node[box=green!12, below=3mm of out] (evi) {Evidence\\+ drivers};
\node[box=green!12, below=3mm of evi] (aud) {Bias\\audit};
\draw[arr] (res) -- (red);
\draw[arr] (jd) -- (req);
\draw[arr] (red) -- (sent);
\draw[arr] (req) -- (sk);
\draw[arr] (sent.east) -- ++(3mm,0) |- ([yshift=2mm]feat.west);
\draw[arr] (sk.east) -- ++(3mm,0) |- ([yshift=-2mm]feat.west);
\draw[arr] (feat) -- (head);
\draw[arr] (head.east) -- ++(2mm,0) |- (out.west);
\draw[arr] (head.east) -- ++(2mm,0) |- (evi.west);
\draw[arr] (head.east) -- ++(2mm,0) |- (aud.west);
\node[lab, above=1mm of red] {removes identity cues};
\node[lab, below=1mm of req] {required / preferred};
\end{tikzpicture}

\caption{EviMatch pipeline. Identity cues are removed before any feature is computed; the same
analysis object yields the feature vector and the evidence shown to the reviewer, so an explanation
cannot describe a signal the model did not use.}
\label{fig:arch}
\end{figure*}

Figure~\ref{fig:arch} gives the overview. Given a resume $R$ and a job description $J$, the system
returns a fit score $S\in[0,100]$, a decision band, and an evidence record.

\subsection{Blind-screening front end}
Redaction is rule based, deterministic and runs before segmentation. It removes e-mail addresses,
telephone numbers, URLs, header names and ``Name:'' fields, personal-detail fields (date of birth,
marital status, nationality, and similar), honorifics, graduation dates, institution names, and
organisation names that carry a gender marker. Gendered pronouns and role nouns are neutralised
(``she'' $\to$ ``they''). Institution names are removed because institutional prestige is a proxy
for socio-economic background; for the same reason the system uses neither college tier nor grade
point average, both of which an earlier version of this project fed to a language model. A redaction
report lists how many items of each class were removed, never their content.

\subsection{Requirement parsing}
A posting mixes duties and qualifications with benefits, legal text and company marketing, which
dilutes any whole-document similarity. The parser splits $J$ into sentence units, assigns each a kind
$\kappa\in\{\text{required},\text{preferred},\text{duty},\text{other}\}$ from section headers and cue
phrases, and drops boilerplate. Weights are fixed constants,
$w_{\text{required}}{=}1$, $w_{\text{duty}}{=}0.75$, $w_{\text{preferred}}{=}0.5$, $w_{\text{other}}{=}0.4$.

\subsection{Requirement-level evidence matching}
Let $q_1,\dots,q_K$ be the requirement units with weights $w_k$, and $s_1,\dots,s_M$ the
sentence units of the redacted resume. With $\mathbf{e}(\cdot)$ an L2-normalised sentence embedding,
\begin{equation}
c_k=\max_{m}\;\mathbf{e}(q_k)^{\!\top}\mathbf{e}(s_m),\qquad
\mathrm{cov}(R,J)=\frac{\sum_k w_k c_k}{\sum_k w_k}.
\label{eq:cov}
\end{equation}
The arg-max sentence is the evidence for requirement $k$. A requirement is shown as \emph{evidenced}
if $c_k\ge0.50$, \emph{partial} if $0.40\le c_k<0.50$ and otherwise \emph{without evidence}. The two
cut-offs were fixed from training-split probes and not tuned afterwards. Because $c_k$ is an embedding
similarity it indicates topical relevance, not demonstrated competence, and the interface says so.

\subsection{Features and fit model}
Each pair yields 22 relational features in four groups: \emph{lexical} (TF-IDF cosine, length-normalised
BM25 with the posting as query, IDF-weighted term overlap); \emph{semantic} (document cosine, the
weighted coverage of Eq.~\eqref{eq:cov}, its required/preferred variants, a lower quartile and minimum
over requirements, and the fractions of requirements above the two cut-offs); \emph{skills}
(IDF- and weight-adjusted coverage of the posting's skills by the resume, Jaccard overlap, and the
number of shared and missing required skills, from a 350-entry transparent lexicon); and
\emph{experience} (estimated years of experience minus the years the posting asks for, an indicator
that the requirement is met, and the same pair for education level). Features that describe only the
resume or only the posting (for example years of experience on its own) are excluded; the next
section explains why.

The head is a three-class model over \{No, Potential, Good Fit\}. We train it without class
re-weighting, so its probabilities $P(c\,|\,x)$ stay calibrated to the base rate, and decide by
dividing by the training prior $\pi_c$:
\begin{align}
\hat c &= \arg\max_c \; P(c\,|\,x)/\pi_c, \label{eq:decide}\\
S &= 100\,\big(P(\text{Good}\,|\,x)+\tfrac12 P(\text{Potential}\,|\,x)\big). \label{eq:score}
\end{align}
Eq.~\eqref{eq:decide} is the Bayes rule for balanced accuracy and maps to the bands Reject, Review and
Shortlist; the middle band is where human review is expected. Candidate heads are a multinomial
logistic regression with an $L_2$ penalty and a histogram gradient-boosted classifier
\cite{pedregosa2011sklearn}; selection uses the one-standard-error rule described in
Section~\ref{sec:setup}.

\subsection{Driver attribution}
The contribution of feature $f$ to a candidate's score is the occlusion effect
\begin{equation}
\Delta_f = S(x)-S(x_{f\leftarrow\mu_f}),
\end{equation}
where $\mu_f$ is the training mean. Occlusion is model agnostic, deterministic and needs no extra library;
effects are also summed per feature group for display. Section~\ref{sec:results} checks it with a
deletion test rather than assuming it is faithful.

\subsection{Counterfactual bias audit}\label{sec:audit}
For a resume $R$ the audit builds variants that change one irrelevant cue and compares the score with a
control. \emph{Name cells}: a name from each of ten race/ethnicity $\times$ gender cells (four names per
cell) is prepended as a header. \emph{Proxy cues}: a final line is appended, always compared with a neutral
line of similar length: an affinity-organisation membership (three groups), stereotyped hobbies (two
codings), a career break for childcare versus for study, and a gender-marked organisation that the
redactor is designed to remove. The statistics are the mean score change with a job-description-cluster
bootstrap interval, the decision-flip rate, and the largest gap between name cells. Name-to-group association
is statistical and imperfect; the lists are fixed in the code so the audit is reproducible.
